%!TeX encoding = UTF-8 Unicode
%!TeX spellcheck = en-UK
% !TeX root = Perrinet26mnesis.tex
% -*- coding: UTF-8; -*-
%%%-----------------------------------------------------------------
\begin{figure}[!t] 
  \centering
      \includegraphics[width=.48\linewidth]{lorenz_chaotic_encoding.pdf}%
    % \hfill
    \includegraphics[width=.48\linewidth]{lorenz_chaotic_decoding.pdf}%
  \caption{%
    We used the equations of the Lorenz attractor to generate a time series of chaotic dynamics. 
    \emph{Left:} We used a simple encoding scheme to convert the continuous dynamics into a sequence of discrete events using a nearest-neighbor one-hot encoding. This allows us to reconstruct the trajectory from spikes. %
    \emph{Right:} 
    After learning, the network is able to autonomously generate the chaotic dynamics of the Lorenz attractor, starting from a short trigger window of input spikes. The network output (red) closely matches the target trajectory (blue) as it encoded by spikes (green) for all three dimensions of the Lorenz attractor.
  }\label{fig:lorenz}
\end{figure}
%%%-----------------------------------------------------------------
